\documentclass[11pt]{article}

\usepackage[margin=1in]{geometry}
\usepackage{enumitem}
\usepackage[hidelinks]{hyperref}

\title{Sharat's Comments}
\author{}
\date{}

\begin{document}

\maketitle

\section*{General Comments}

\begin{itemize}[leftmargin=*, itemsep=0.5em]
    \item Title must be changed. Information gaps gives very little information about what we are doing. It makes it sound like we are working with incorrect/incomplete information.   

    Some okayish candidates: Power of Adaptivity in the Stochastic Traveling Salesperson Problem; Adaptive Algorithms for the Stochastic Traveling Salesperson Problem. 

    We should come up with a name that can capture what you accomplished in your thesis while not being too generic or too specific.

\end{itemize}

\section*{Detailed Comments}

\begin{itemize}[leftmargin=*, itemsep=0.5em]
    \item \textbf{Page a, line b:}
    
\end{itemize}

\end{document}


\documentclass[11pt]{article}

\usepackage[T1]{fontenc}
\usepackage{lmodern}
\usepackage[margin=1in]{geometry}
\usepackage{amsmath,amssymb,amsthm,mathtools}

\newtheorem{lemma}{Lemma}

\newcommand{\eps}{\varepsilon}
\newcommand{\cost}{\operatorname{cost}}
\newcommand{\loc}{\mathrm{loc}}

\begin{document}

\begin{lemma}[Pairwise truncation bound]
\label{lem:long-gap-truncation}
Let $0<\eps\le 1$ and $0<\eta\le \eps^2/16$.  For any a~priori
tour~$T$, we have
\[
    \cost^{\loc}_{\eta}(T)
    \le
    \cost(T)
    \le
    (1+\eps)\cost^{\loc}_{\eta}(T).
\]
\end{lemma}

\begin{proof}
The first inequality is immediate because the truncated cost omits only
nonnegative terms.  We prove the second inequality.

For $0\le i<j\le n+1$, let
\[
    w_{ij}\coloneqq p_i p_j Q(i,j),
\]
and write
\[
    C\coloneqq\cost(T),
    \qquad
    C_\eta\coloneqq\cost^{\loc}_{\eta}(T),
    \qquad
    D_\eta\coloneqq C-C_\eta.
\]
Thus, $D_\eta$ is the total contribution of the discarded pairs, namely
those pairs $(i,j)$ for which $Q(i,j)\le\eta$.

Fix a discarded pair $(i,j)$.  Independently sample the activation states of
the vertices strictly between $v_i$ and $v_j$, conditional on at least one of
them being active.  This conditioning event has probability
\[
    1-Q(i,j)\ge 1-\eta.
\]
Together with $v_i$ and $v_j$, the sampled active vertices form a path in the
order in which they occur in~$T$.  Denote its edge set by $P_{ij}$.  By the
triangle inequality,
\[
    d_{ij}
    \le
    \sum_{(a,b)\in P_{ij}} d_{ab}.
\]
We use this inequality to charge $w_{ij}d_{ij}$ to the edges of the random
path~$P_{ij}$.

Fix a target pair $(a,b)$ with $w_{ab}>0$, and let $\Gamma_{ab}$ be the
expected coefficient charged to $d_{ab}$.  A source pair $(i,j)$ can charge
$(a,b)$ only if $i\le a<b\le j$.  If $(i,j)\ne(a,b)$, then, before
conditioning, the probability that $a$ and $b$ are consecutive on the sampled
path is
\[
    Q(a,b)
    p_a^{\mathbf 1[i<a]}
    p_b^{\mathbf 1[b<j]}.
\]
Conditioning increases this probability by at most a factor $1/(1-\eta)$.
After dividing by $w_{ab}$ and separating the three possible strict-containment
patterns, we obtain
\begin{align}
\frac{\Gamma_{ab}}{w_{ab}}
\le \frac{1}{1-\eta}
\biggl(&
    \sum_{\substack{j>b\\Q(a,j)\le\eta}} p_jQ(a,j)
    +
    \sum_{\substack{i<a\\Q(i,b)\le\eta}} p_iQ(i,b)
    \notag\\
    &+
    \sum_{\substack{i<a,\ j>b\\Q(i,j)\le\eta}}
        p_ip_jQ(i,j)
\biggr).
\label{eq:congestion}
\end{align}

The first two sums in~\eqref{eq:congestion} are each at most~$\eta$.  Indeed,
$Q(a,j)$ is nonincreasing in~$j$, so the indices satisfying $Q(a,j)\le\eta$
form a suffix, and telescoping gives
\[
    \sum_{\substack{j>b\\Q(a,j)\le\eta}}p_jQ(a,j)
    \le \eta.
\]
The second sum is bounded symmetrically.

It remains to bound the third sum in~\eqref{eq:congestion}.  Set
\[
    H\coloneqq (1-p_a)Q(a,b)(1-p_b),
    \qquad
    x_i\coloneqq Q(i,a),
    \qquad
    y_j\coloneqq Q(b,j).
\]
For $i<a<b<j$, we have $Q(i,j)=Hx_iy_j$.  Moreover, telescoping gives
\[
    \sum_{i<a}p_ix_i=1,
    \qquad
    \sum_{j>b}p_jy_j=1.
\]
If either index set is empty, the third sum is zero.  If $H\le\eta$, it is at
most~$H$ and hence at most~$\eta$.  Suppose, therefore, that $H>\eta$, and put
$t\coloneqq\eta/H$.  For every $z>0$, telescoping on the right gives
\[
    \sum_{\substack{j>b\\y_j\le z}}p_jy_j
    \le \min\{1,z\}.
\]
Terms with $x_i=0$ vanish.  Consequently,
\begin{align*}
    \sum_{\substack{i<a,\ j>b\\Q(i,j)\le\eta}}
        p_ip_jQ(i,j)
    &=
    H\sum_{\substack{i<a\\x_i>0}}p_ix_i
        \sum_{\substack{j>b\\y_j\le t/x_i}}p_jy_j\\
    &\le
    H\sum_{i<a}p_i\min\{x_i,t\}.
\end{align*}

We use the following elementary survival-product estimate:
\begin{equation}
    \sum_{i<a}p_i\min\{Q(i,a),t\}
    \le
    t\left(1+\ln\frac1t\right)
    \qquad (0<t<1).
\label{eq:survival-product}
\end{equation}
To verify it, recall that $x_i=Q(i,a)$ is nondecreasing in~$i$ and that
\[
    x_{i-1}=(1-p_i)x_i.
\]
If $x_0>t$, then
\[
    \sum_{i<a}p_i\min\{x_i,t\}
    =t\sum_{i<a}p_i
    \le t\left(1+\ln\frac1{x_0}\right)
    \le t\left(1+\ln\frac1t\right).
\]
Otherwise, let $\ell$ be the largest index for which $x_\ell\le t$, and put
$x\coloneqq x_{\ell+1}$ and $z\coloneqq x_\ell$.  Since
$z=(1-p_{\ell+1})x$, telescoping and the inequality
$u\le-\ln(1-u)$ give
\begin{align*}
    \sum_{i<a}p_i\min\{x_i,t\}
    &\le
    z+t\left(1-\frac{z}{x}+\ln\frac1x\right)\\
    &\le
    t\left(2-\frac{t}{x}+\ln\frac1x\right)\\
    &\le
    t\left(1+\ln\frac1t\right).
\end{align*}
The last inequality is equivalent to
$1-t/x+\ln(t/x)\le0$, which follows from $\ln u\le u-1$.
This proves~\eqref{eq:survival-product}.

Applying~\eqref{eq:survival-product}, the third sum in
\eqref{eq:congestion} is at most
\[
    Ht\left(1+\ln\frac1t\right)
    =
    \eta\left(1+\ln\frac{H}{\eta}\right)
    \le
    \eta\left(1+\ln\frac1\eta\right).
\]
It follows from~\eqref{eq:congestion} that every target pair has congestion at
most
\[
    \beta(\eta)
    \coloneqq
    \frac{\eta\bigl(3+\ln(1/\eta)\bigr)}{1-\eta}.
\]
Pairs with $w_{ab}=0$ receive no charge, so summing over all target pairs gives
\[
    D_\eta\le \beta(\eta)C.
\]
Therefore,
\begin{equation}
    C_\eta
    \ge
    \bigl(1-\beta(\eta)\bigr)C.
\label{eq:retained-fraction}
\end{equation}

The function $\beta$ is increasing on $(0,1)$.  Since
$\eta\le\eps^2/16\le1/16$, we have
\begin{align*}
    \beta(\eta)
    &\le
    \beta\left(\frac{\eps^2}{16}\right)\\
    &\le
    \frac{\eps^2}{15}
    \left(3+\ln 16+2\ln\frac1\eps\right)\\
    &\le
    \frac{\eps^2}{15}
    \left(4+\frac{2}{\eps}\right)\\
    &\le
    \frac{2\eps}{5}
    \le
    \frac{\eps}{1+\eps}.
\end{align*}
Here we used $\ln(1/\eps)\le 1/\eps-1$, $\ln 16<3$, and
$0<\eps\le1$.  Substituting this bound into~\eqref{eq:retained-fraction}
gives
\[
    C_\eta
    \ge
    \frac{C}{1+\eps},
\]
which is the desired second inequality.
\end{proof}

\end{document}